% !TEX root = ../main.tex

\chapter{绪论}

\section{研究背景与意义}

随着具身智能、移动机器人和深度学习技术的快速发展，机器人正在从结构化、封闭式的实验室环境走向消防巡检、室内导览、园区巡逻、地下管廊检查和灾后排查等真实应用场景。在这些场景中，机器人不仅要感知周围环境，还要在地图不完备、光照变化明显、有障碍物、通道空间受限以及定位先验较弱的条件下完成自主移动。导航因此不再只是移动机器人“从一点到另一点”的路径计算问题，而是机器人完成巡检、导览、搬运、搜救和服务任务之前必须具备的基础能力。没有稳定的导航能力，感知结果难以转化为有效行动，任务规划也无法落到具体空间执行上。因此，导航逐渐成为连接视觉感知、任务理解、局部决策和底层执行的核心环节\cite{pei2026embodied_nav}。

从应用需求来看，让机器人进入上述场景具有直接的现实意义。消防巡检、灾后排查和地下管廊检查等任务往往伴随烟尘、高温、有毒气体、坍塌风险或长时间重复劳动，完全依赖人工会带来较高安全成本和人力消耗；室内导览、园区巡逻、楼宇安防和场馆服务等任务虽然危险性较低，但也对连续运行、路径复现和场景适应能力提出了要求。与人相比，机器人可以在危险或重复环境中保持稳定执行，并通过视觉传感器持续记录环境状态，因此适合承担移动感知与目标到达类任务。然而，不同机器人平台对环境的适应能力和工程部署条件存在明显差异，因此有必要进一步明确本文所关注的机器人形态。

从机器人平台类型来看，不同形态的移动机器人具有不同的适用边界。人形机器人在结构上更接近人类，理论上能够适应面向人设计的环境，并完成一定的人机交互任务，但其系统结构复杂、控制难度高，稳定性和工程部署成本仍然较高。轮式底盘机器人结构简单、控制成熟、能耗较低，在平整地面和规则室内环境中具有较好的应用基础，但其运动能力较依赖地面条件，面对台阶、障碍物、碎石、坡道或非结构化地形时通过性受到限制。相比之下，四足机器人在机动性、稳定性和环境适应性之间具有较好的平衡：其足式运动方式能够跨越一定高度的障碍物，适应复杂地面和狭窄空间，同时相比人形机器人具有更低的结构复杂度和更高的工程可部署性。因此，在灾害现场、工业巡检、园区巡逻和室内外混合环境等任务中，四足机器人是一类具有代表性的具身导航平台\cite{yang2019quadruped_review,hwangbo2019anymal,lee2020learning}。

基于上述分析，如何使四足机器人根据简单、直观的目标表达，在真实环境中自主到达指定位置，是具身导航研究面向工程落地时必须回答的问题。传统移动机器人导航通常依赖全局坐标、度量地图和预设航点。这一路线在仓储物流、工厂搬运、医院配送、商场清洁和室内导览等平整、规则环境中已经形成较成熟的系统链路，轮式平台也在这些场景中承担了重复运输、固定路线巡检和服务引导等角色。然而，在临时巡检、陌生室内空间或弱 GPS 场景中，精确地图往往难以及时获取，普通用户也不容易用坐标描述目标位置。在真实应用中，操作者往往更容易通过目标图像、路线图像或关键场景记忆表达“想让机器人到达哪里”。

视觉目标导航为这一问题提供了更自然的任务表达方式：操作者只需提供目标位置的第一视角图像，或由若干关键图像组成的拓扑地图，机器人便可根据当前图像与目标图像之间的关系生成局部运动。这种“看见某处并走过去”的形式更贴近人的空间记忆方式，也更适合导览复现、巡检路线记录和弱先验环境中的目标到达任务\cite{zhu2017targetdriven,shah2021ving}。但是，对于四足平台而言，并不是将轮式导航算法直接更换载体即可使用。四足机器人的步态运动会引入相机抖动和机体俯仰，速度指令与真实位移之间也会存在非线性偏差，高层视觉策略还必须与底层运动主机、通信协议和安全停机机制形成稳定闭环。

近年来，通用视觉导航模型和生成式策略的发展为上述问题提供了新的技术基础。通用导航模型（General Navigation Model, GNM）和视觉导航 Transformer（Visual Navigation Transformer, ViNT）将历史图像观察、目标条件和局部航点预测纳入统一模型\cite{gnm2022,vint2023}；目标掩码扩散导航模型（Navigation with Goal Masked Diffusion, NoMaD）进一步引入动作扩散策略，通过条件去噪采样生成多条局部轨迹候选，并利用目标掩码机制（goal mask）统一目标导航与探索行为\cite{nomad2023,chi2025diffusion}。与直接回归单一路径相比，扩散式轨迹生成能够描述多模态决策分布，更适合走廊分叉、障碍绕行和目标不确定等情形。然而，现有公开验证仍多集中在轮式平台或离线数据集上，面向四足机器人真实部署时的实时推理、视觉扰动、运动接口适配和安全闭环问题仍需要系统研究。

基于上述背景，本文以 NoMaD 为高层视觉导航模型，以云深处 Lite3 四足机器人为部署对象，研究在室内狭窄、弱先验、无全局精确地图的环境中，如何将扩散式视觉目标导航模型部署到资源受限且运动扰动明显的四足平台上，使其能够根据真实相机图像与目标图像形成可实时运行、可仿真验证、可真机执行的闭环导航系统。

\begin{figure}[htbp]
    \centering
    \includegraphics[width=0.5\linewidth]{intro-background-placeholder.png}
    \caption{四足机器人视觉目标导航目标场景示意图}
    \label{fig:intro-background-placeholder}
\end{figure}

围绕这一问题，本文的研究意义主要体现在三个层面。首先，在应用层面，本文将视觉目标导航与四足平台结合，使导航任务不再停留于离线图像匹配或仿真指标，而是面向室内巡检、导览和受限空间通行等具体场景。其次，在算法层面，本文关注动作扩散策略在边缘设备上的推理效率和轨迹质量，通过去噪扩散隐式模型（Denoising Diffusion Implicit Model, DDIM）采样加速、无分类器引导（Classifier-Free Guidance, CFG）、测试时扩展（Test-Time Scaling, TTS）候选筛选和视觉编码器替换等实验，分析扩散式导航模型在真实部署中的速度、质量和稳定性权衡。最后，在系统层面，本文构建从统一推理模块、状态机、导航主机、中层控制映射到 Lite3 UDP 运动接口的分层闭环架构，使高层视觉导航策略能够接入 MuJoCo 仿真环境和 Jetson Orin 真机平台，为后续四足视觉导航研究提供可复现的工程基础。

\section{国内外研究现状}

四足机器人视觉目标导航涉及传统导航与视觉 SLAM、学习型视觉导航、通用视觉导航模型、动作扩散策略以及足式机器人真实部署等多个方向。现有研究已经分别在地图构建、目标表达、轨迹生成、跨平台泛化和四足运动控制方面取得了重要进展，但这些方向往往各有侧重：传统导航强调定位建图和规划可靠性，学习型视觉导航强调从图像到动作的策略学习，通用导航模型强调跨场景和跨机器人泛化，四足机器人研究则更多关注底层运动控制和复杂地形通过。本文关注的问题位于这些方向的交叉处：在弱先验视觉目标导航任务中，如何把生成式视觉导航策略与四足机器人的真实运动执行、边缘计算和安全闭环结合起来，形成可实时运行、可安全执行、可从仿真验证到真机部署的系统。

\subsection{传统导航、视觉 SLAM 与学习型导航}

传统机器人导航通常以状态估计、地图构建、路径规划和运动控制为核心。Thrun 等\cite{thrun2005probabilistic}建立了概率机器人学框架，使机器人能够在不确定感知和运动条件下进行定位、建图与决策；该工作奠定了传统导航的理论基础，做到了对传感器噪声、运动误差和环境不确定性的概率化建模，但它并不直接解决由目标图像驱动的视觉导航问题，也不直接涉及学习型视觉导航策略在四足平台上的闭环实现。Hess 等\cite{cartographer2016}提出 Cartographer，实现了二维激光 SLAM 中实时回环检测和子图优化，在激光雷达配置充分的轮式平台上具有较强工程可靠性；但其重点仍是度量地图构建，并不面向纯视觉单目 RGB 目标图像导航问题。Campos 等\cite{orbslam3_2021}提出 ORB-SLAM3，在视觉、视觉惯性和多地图 SLAM 中取得了较高定位精度；该方法做到了成熟的视觉几何定位，但性能依赖特征质量、标定精度和几何一致性，不以图像目标到达和动作策略生成为核心。

国内研究也对 SLAM、具身导航和机器人导航系统进行了系统梳理。毛军等\cite{mao2022slam_review}总结了惯性、视觉和激光雷达 SLAM 的技术演进，做到了从传感器融合和定位建图角度梳理导航基础能力；但其关注点主要是环境表示与位姿估计，而不是目标图像条件下的端到端导航决策。裴凌和熊超然\cite{pei2026embodied_nav}从具身导航的概念、方法和挑战出发，指出导航研究正从几何到达转向任务驱动、语义理解与自主决策；该综述为本文研究提供了问题背景，但并未给出四足机器人平台上的视觉导航部署链路。因此，传统导航和 SLAM 工作为机器人自主移动提供了重要基础，但仍需要结合学习型视觉策略和真实平台系统实现，才能支撑本文所关注的弱先验四足视觉目标导航任务。

深度学习方法进一步推动了视觉导航从几何建模向策略学习转变。Zhu 等\cite{zhu2017targetdriven}提出目标驱动视觉导航方法，利用深度强化学习建立当前图像、目标图像与动作之间的映射；该工作证明了图像目标可以作为导航任务输入，但主要在仿真环境中验证，距离真实四足机器人闭环部署仍有差距。Anderson 等\cite{anderson2018evaluation}讨论具身导航智能体的评价问题，推动了导航成功率、路径效率等指标的规范化；该工作做到了评价标准层面的统一，但评价指标本身不能替代真实平台上的感知、通信、控制和安全闭环设计。Wijmans 等\cite{ddppo2020}通过大规模分布式强化学习训练点目标导航智能体，在仿真中取得了很高性能；但点目标导航依赖相对坐标目标，而不是视觉目标导航。

\subsection{视觉目标导航与拓扑记忆}

视觉目标导航通过目标图像、目标物体或语言指令描述机器人需要到达的位置。相较于精确坐标，目标图像更适合弱先验环境。陈铂垒等\cite{chen2025object_nav_review}综述了物体目标导航中的语义地图、强化学习和大模型辅助方法，说明视觉语义目标已经成为具身导航的重要方向。司马双霖等\cite{sima2023vln_review}梳理了视觉语言导航中语言、视觉和路径规划的结合方式，做到了对多模态导航任务的系统总结；但语言目标存在语义歧义和解析成本，本文选择采用目标图像和拓扑图像节点作为更直接的任务表达。

拓扑记忆方法通过关键图像节点组织长程路线，为弱地图先验下的视觉导航提供了重要思路。Savinov 等\cite{savinov2018sptm}提出半参数拓扑记忆 SPTM，利用视觉相似度和局部策略沿记忆图前进；该方法做到了用稀疏视觉节点替代完整度量地图，但局部控制策略和节点质量仍会影响实际执行效果。Chaplot 等\cite{chaplot2020nts}提出神经拓扑 SLAM，将学习表示与拓扑规划结合，在复杂室内场景中构建拓扑图并规划路径；该工作增强了拓扑规划能力，但系统复杂度较高，也并非面向四足平台。Shah 等\cite{shah2021ving}提出 ViNG，通过视觉路标和目标图像实现开放世界视觉目标导航，并在真实环境中展示了长程导航能力；该工作做到了目标图像导航和真实机器人验证，但主要面向轮式平台，并未重点讨论四足机器人步态扰动、航点到速度指令映射、安全状态机和真机闭环频率等问题。

因此，视觉目标导航和拓扑记忆已经证明了“用图像表达目标”的可行性。本文并不把局部航点或拓扑节点本身作为研究问题，而是将其作为解决视觉目标导航的手段：系统在拓扑图像序列中选择实时子目标，再由基于目标掩码的扩散策略生成局部动作，并进一步映射为四足机器人可执行的速度指令。

\subsection{动作扩散策略与通用视觉导航模型}

扩散模型为多模态轨迹生成提供了新的建模工具。Ho 等\cite{ddpm2020}提出 DDPM，通过逐步加噪和反向去噪学习复杂数据分布，做到了高质量生成建模，但迭代采样会带来较高推理开销。Song 等\cite{ddim2020}提出 DDIM，在不重新训练模型的情况下减少反向采样步数，做到了扩散模型推理加速。Ho 和 Salimans\cite{cfg2022}提出无分类器引导，为条件生成提供了可调的条件增强方式；该方法提升了条件一致性。Chi 等\cite{chi2025diffusion}提出 Diffusion Policy，将扩散模型用于机器人视觉运动控制，做到了以动作序列生成方式表达多模态行为；但该工作主要面向操作控制任务，并不直接解决机器人导航问题。

在视觉导航领域，GNM、ViNT 和 NoMaD 形成了通用导航模型的重要脉络。Shah 等\cite{gnm2022}提出 GNM，尝试用统一模型驱动不同机器人平台，证明了跨机器人视觉导航模型的可行性；但 GNM 主要采用轨迹回归形式，没有引入动作扩散生成机制，也没有统一目标导航和无目标探索。ViNT\cite{vint2023}进一步采用 Transformer 融合历史观察和目标图像，做到了更强的时序视觉表征和目标条件建模；但其重点仍是视觉导航基础模型，不以扩散采样优化和真实部署为核心。Sridhar 等\cite{nomad2023}提出 NoMaD，将目标掩码机制与动作扩散策略结合，在同一模型中统一目标导航和无目标探索，做到了用统一扩散策略生成多模态局部轨迹。NoMaD 是本文复现、优化和部署的直接基础；但原工作主要展示通用导航模型能力，公开实验主线更偏向轮式移动平台和通用导航评估，对足式平台上的视觉扰动、速度接口、安全状态和端侧算力约束讨论较少。

近期扩散式视觉导航研究进一步提升了轨迹生成效率和泛化能力。NaviBridger\cite{navibridger2025}提出基于去噪扩散桥的视觉导航框架，从有信息的先验动作而非纯高斯噪声开始生成动作序列，做到了在仿真和真实室内外场景中提升推理效率和动作质量；其已有真实机器人实验，但研究重点仍是有目标条件下的扩散桥动作生成，对目标导航与无目标探索的统一系统，以及四足平台闭环部署效率讨论较少。NavDP\cite{navdp2025}提出基于特权信息指导的导航扩散策略，利用大规模仿真数据训练，并强调跨机器人形态的零样本泛化 sim2real，在轮式、四足和人形平台上展示了较强泛化能力。但其核心贡献是大规模仿真训练与跨形态泛化，足式平台更多作为跨形态验证对象；对单目图像目标/拓扑目标驱动的四足闭环系统、速度接口适配和闭环频率评估并非其重点。SIDP\cite{sidp2026}提出自模仿扩散策略，通过奖励引导的自模仿减少对大量采样和后处理筛选的依赖，报告了在 Jetson Orin Nano 上相较 NavDP 更快的纯推理时间，但其重点是路径推理延迟和模型效率，系统层面的真机闭环频率与完整执行链路不是其主要评价对象。因此，本文与这些工作面向的具体问题不同，使用的方法和评价层级也不同：本文更关注这类扩散式视觉导航方法在具体四足机器人平台上的系统部署、闭环效率和 sim2real 实验分析。

\subsection{四足机器人视觉导航与真实部署}

四足机器人研究长期关注运动能力、复杂地形通过和感知运动控制。杨钧杰等\cite{yang2019quadruped_review}综述了四足机器人的机构设计、步态规划和控制方法，做到了对国内外四足机器人研究脉络的整理；但该综述主要面向运动控制和平台发展，对扩散策略讨论较少。Hwangbo 等\cite{hwangbo2019anymal}利用强化学习使 ANYmal 获得敏捷动态运动能力，Lee 等\cite{lee2020learning}进一步研究四足机器人在复杂地形上的学习型运动控制；这些工作做到了真实四足平台上的鲁棒运动控制，但没有把目标图像导航作为高层导航问题来处理。Miki 等\cite{miki2022perceptive}通过感知与本体信息融合提升四足机器人野外通过能力，说明视觉和地形感知对足式机器人具有重要价值；但其关注点是复杂地形中的感知运动控制，而不是导航的闭环部署。

近年来，四足视觉导航和多模态四足机器人任务开始受到更多关注。Ren 等\cite{ren2023hvn}提出分层视觉导航系统，将视觉导航与落足点适应学习结合，用于提升四足机器人在复杂环境中的导航能力；该工作做到了四足平台上的视觉导航与落足适应结合，但其重点是视觉导航与落足控制协同，对图像目标条件下的生成式轨迹分布和边缘部署效率讨论较少。Kareer 等\cite{kareer2023vinl}提出 ViNL，将视觉导航与越障运动控制结合，使四足机器人能够在室内障碍环境中执行导航任务；该工作证明了四足视觉导航的可行性，但重点是视觉导航和运动控制协同训练，并非面向目标图像/拓扑图像驱动的扩散式导航部署。DiPPeST\cite{dippest2024}提出面向四足机器人的图像与目标条件扩散轨迹规划器，并通过视觉伺服框架进行了真实四足机器人执行验证，说明扩散轨迹生成确实已经被用于四足机器人路径规划；但它更偏向扩散路径规划器和视觉伺服框架本身，对目标图像或拓扑节点驱动的导航闭环、边缘设备效率和系统失败模式分析仍不充分。QUAR-VLA\cite{quarvla2024}提出四足机器人的视觉-语言-动作模型和多任务数据集，将视觉、语言和动作结合起来生成可执行动作；该工作拓展了四足机器人多模态任务能力，但其重点是 VLA 范式和语言条件任务，不是面向弱先验室内场景的单目图像目标导航部署。

综上，现有研究已经证明了四足机器人可以进行视觉导航，也证明了扩散策略可以用于导航和轨迹生成。从真实部署角度看，仍有三个方面值得进一步系统化讨论：第一，视觉目标导航与四足机器人研究多沿着各自路线推进，前者更关注目标表达、策略学习和跨场景泛化，后者更关注运动控制、落足适应和复杂地形通过，弱先验图像目标任务下高层视觉决策与足式执行接口之间的耦合关系仍需要结合具体平台分析；第二，扩散式导航研究多强调模型结构、训练规模、采样策略或纯推理时间，边缘计算部署、采样加速、候选筛选、速度指令映射和安全状态管理在同一真实闭环链路中的综合评估仍有进一步展开空间；第三，已有 sim2real 工作通常更重视跨平台泛化能力，针对具体四足平台从仿真到真机的闭环频率、任务成功率、失败模式和工程边界的系统化分析仍可补充。本文正是围绕这些问题展开，构建面向四足机器人的扩散式视觉目标导航闭环系统，完成目标导航与探索统一推理链路的部署优化与平台适配，并给出从仿真到真实四足机器人的系统验证与分析。

\section{本文主要贡献}

针对四足机器人在弱先验环境中的视觉目标导航部署问题，本文并不试图提出规模更大的通用导航基础模型，而是面向真实四足机器人平台，围绕扩散式视觉目标导航策略的闭环部署、推理优化、平台适配和 sim2real 验证开展研究。本文主要贡献如下：

\begin{enumerate}
    \item 构建了面向四足机器人的扩散式视觉目标导航闭环系统。本文围绕四足机器人视觉目标导航任务，构建了从视觉观测、目标图像输入、扩散式动作生成到底层运动控制的完整闭环系统。系统以单目 RGB 图像和目标图像为输入，通过视觉编码、距离预测和扩散策略生成局部导航动作，并进一步转换为四足机器人可执行的速度控制指令，实现了扩散式视觉导航方法在真实四足机器人平台上的系统化落地。本文以云深处 Lite3 作为具体实验平台完成验证。
    \item 完成了基于目标掩码的扩散策略统一推理链路的部署优化与四足平台适配。针对基于目标掩码的扩散式视觉导航策略在真实机器人部署中存在的推理延迟、动作接口不匹配和任务状态切换等问题，本文构建了统一推理链路，并完成了面向四足机器人平台的适配。
    \item 给出了扩散式视觉导航策略从仿真到真机的系统验证与分析。本文在仿真环境和真实四足机器人平台上对所构建系统进行了验证，分析了不同扩散采样方式和推理配置对导航成功率、推理效率与闭环运行频率的影响。实验结果表明，经过部署优化后的系统能够在真实场景中完成视觉目标导航与探索任务，并在真机闭环中达到最高 6.80\,Hz 的运行频率。同时，本文进一步分析了方法在光照变化、视角差异、目标图像匹配误差和底层运动控制误差等因素下的局限性，为扩散式视觉导航策略的 sim2real 部署提供了实验依据。
\end{enumerate}

\section{论文章节安排}

本文整体内容按照“问题提出—统一推理模块构建—系统闭环验证—真机部署—总结展望”的逻辑展开，共分为五章。各章节之间的逻辑结构关系如图~\ref{fig:chapter-logic} 所示。第一章首先明确研究场景与问题边界，第二章围绕高层视觉导航策略构建基于动作扩散策略的统一视觉目标导航推理模块，第三章进一步将算法接入四足机器人闭环系统并完成仿真验证，第四章面向真实硬件平台完成部署与实验，第五章在前文结果基础上总结全文并讨论后续研究方向。

\begin{figure}[htbp]
    \centering
    \includegraphics[width=0.8\linewidth]{structure.png}
    \caption{本文各章节逻辑结构关系}
    \label{fig:chapter-logic}
\end{figure}

第一章为绪论，介绍四足机器人视觉目标导航的研究背景与意义，梳理传统导航、视觉目标导航、动作扩散策略和四足真实部署等方向的国内外研究现状，并明确本文的研究问题、主要贡献和章节结构。

第二章围绕基于动作扩散策略的统一视觉目标导航推理模块展开，首先给出视觉目标导航、拓扑地图、航点、扩散模型和 NoMaD 机制等必要背景，随后介绍数据集、实验设备、基线推理验证、DDIM、CFG、TTS 与视觉编码器替换等推理模块实验。

第三章介绍基于分层架构与 MuJoCo 仿真的四足机器人视觉导航闭环系统设计与验证，说明轮式平台与足式平台的部署差异、Lite3 平台基础、统一推理模块、状态机、导航主机、中层控制映射、MuJoCo 仿真平台和闭环实验结果。

第四章给出 Jetson Orin 与 Lite3 真机部署流程，说明真机系统双链路、Orin 部署架构、真实拓扑地图构建、滑动窗口导航、速度指令下发、安全机制和真实环境实验结果。

第五章总结全文工作，归纳本文在算法优化、系统集成和真机验证方面的成果，并讨论后续可进一步拓展的研究方向。

\section{本章小结}

本章首先从具身智能和真实应用需求出发，说明了四足机器人视觉目标导航在消防巡检、室内导览和受限空间通行等场景中的研究价值，并进一步分析了视觉目标表达、四足平台形态和动作扩散策略带来的机遇与挑战。随后，本章按照导航建模、目标表达、扩散策略和四足部署的技术逻辑，对国内外相关研究进行了梳理，指出现有工作在真实四足平台上的扩散式视觉目标导航部署方面仍存在研究空间。最后，本章明确了本文面向 Lite3 四足机器人的研究内容和章节安排，为后续章节展开奠定基础。
